\section{Multi-GPU Oceananigans on Nibi (Quick Start)} 

This appendix describes the minimal setup required to run Oceananigans benchmarks on Nibi using H100 GPUs and CUDA-aware MPI. 

\subsection{Login} 

\begin{verbatim} 
    ssh fpoulin@nibi.alliancecan.ca 
\end{verbatim} 

\subsection{Load Modules} 

Start from a clean environment: 

\begin{verbatim} 
    module purge 
    module load StdEnv/2023 
    module load gcc/12.3 
    module load openmpi/4.1.5 
    module load cuda/12.6 
    module load julia/1.10.10 
\end{verbatim} 

Verify: 

\begin{verbatim} 
    module list julia --version 
\end{verbatim} 

Expected: 

\begin{verbatim} 
    julia version 1.10.10 
\end{verbatim} 

\subsection{Move the Julia Depot to Scratch} 

To avoid home-directory quota issues: 

\begin{verbatim} 
    export JULIA_DEPOT_PATH=$SCRATCH/julia_depot 
    
    mkdir -p $JULIA_DEPOT_PATH 
\end{verbatim} 

Verify: 

\begin{verbatim} 
    echo $JULIA_DEPOT_PATH 
\end{verbatim} 

\subsection{Clone Oceananigans} 

Work in scratch rather than home: 

\begin{verbatim} 
    cd $SCRATCH 
    mkdir -p software 
    cd software 
    git clone git@github.com:CliMA/Oceananigans.jl.git 
    cd Oceananigans.jl 
\end{verbatim} 

\subsection{Instantiate Oceananigans} 

\begin{verbatim} 
    julia --project=. 
\end{verbatim} 

Inside Julia: 

\begin{verbatim} 
    using Pkg 
    Pkg.instantiate() 
\end{verbatim} 

Exit Julia. 

\subsection{Instantiate the Benchmarking Environment} 

\begin{verbatim} 
    cd benchmarking 
    julia --project=. 
\end{verbatim} 

Inside Julia: 

\begin{verbatim} 
    using Pkg 
    Pkg.instantiate() 
\end{verbatim} 

Exit Julia. 

\subsection{Request an Interactive H100 Node} 

Do not request only a GPU. Julia may be killed by the memory cgroup if too little system RAM is requested. 

Recommended: 

\begin{verbatim} 
    salloc \ 
    --account=def-fpoulin \ 
    --partition=gpubase_interac \ 
    --gpus=h100:1 \ 
    --cpus-per-task=8 \ 
    --mem=32G \ 
    --time=01:00:00 
\end{verbatim} 

Larger memory requests (e.g. 64G) are recommended for benchmarking. 

\subsection{Verify the GPU} 

\begin{verbatim} 
    hostname 
    nvidia-smi 
\end{verbatim} 

Example: 

\begin{verbatim} 
    NVIDIA H100 80GB HBM3 
\end{verbatim} 

\subsection{Verify CUDA} 

Start Julia: 
\begin{verbatim} 
    julia --project=. 
\end{verbatim} 

Then run: 

\begin{verbatim} 
    using CUDA 
    CUDA.functional() 
    CUDA.versioninfo() 
\end{verbatim} 

Expected: 

\begin{verbatim} 
    true 
\end{verbatim} 

\subsection{Configure MPI.jl to Use System OpenMPI} 

By default MPI.jl may install MPICH\_jll. Configure MPI.jl to use the CUDA-aware OpenMPI installation provided by Nibi. 

Inside Julia: 

\begin{verbatim} 
    using MPIPreferences 
    
    MPIPreferences.use_system_binary( 
        library_names=["libmpi"], 
        extra_paths=[ 
            "/cvmfs/soft.computecanada.ca/easybuild/software/2023/" * 
            "x86-64-v4/Compiler/gcc12/openmpi/4.1.5/lib" 
            ]) 
\end{verbatim} 

Restart Julia. 

Verify: 

\begin{verbatim} 
    using MPI 
    MPI.versioninfo() 
\end{verbatim} 

Expected: 

\begin{verbatim} 
    binary: system 
    abi: OpenMPI 
    Open MPI v4.1.5 
\end{verbatim} 

\subsection{Verify CUDA-Aware MPI} 

\begin{verbatim} 
    using MPI 
    using CUDA 
    MPI.has_cuda() 
\end{verbatim} 

Expected: 

\begin{verbatim} 
    true 
\end{verbatim} 

\subsection{Simple MPI+CUDA Test} 

Create:

\begin{verbatim} 
    using MPI 
    using CUDA MPI.Init()

    rank = MPI.Comm_rank(MPI.COMM_WORLD) 

    println("Rank = ", rank) 
    println("Device = ", CUDA.device()) 
    println("CUDA aware = ", MPI.has_cuda()) 
    
    MPI.Finalize() 
\end{verbatim} 

Run: 

\begin{verbatim} 
    mpiexec -n 1 julia --project=. test_gpu_mpi.jl 
\end{verbatim} 

Expected: 

\begin{verbatim} 
    Rank = 0 
    Device = CuDevice(0) 
    CUDA aware = true 
\end{verbatim} 

\subsection{Run a Benchmark} 

Move into the benchmarking directory: 

\begin{verbatim} 
    cd benchmarking 
\end{verbatim} 

Example benchmark: 

\begin{verbatim} 
    julia --project=. run_benchmarks.jl \ 
    --size=720x360x100 \ 
    --device=GPU 
\end{verbatim} 

This benchmark uses approximately 

\begin{verbatim} 
    6.5 GiB 
\end{verbatim} 

of GPU memory on an H100 and achieves approximately 

\begin{verbatim} 
    1.94e8 grid points/second 
\end{verbatim} 

for the Earth-Ocean benchmark. 

\subsection{Notes} 

\begin{itemize} 
    \item Use scratch for all Oceananigans work. 
    \item Use scratch for the Julia depot. 
    \item Nibi provides CUDA-aware OpenMPI; rebuilding MPI is not necessary. 
    \item Always launch MPI-enabled scripts with: 
    \begin{verbatim} 
        mpiexec -n N julia --project=. script.jl 
    \end{verbatim} 
    \item Always specify memory in GPU allocations. 
\end{itemize}
